\begin{table*}[htbp]
\centering
\scriptsize
\caption{Decomposition of the style index into stylistic and structural parts, two series}
\label{tab:s14}
\begin{tabular}{>{\RaggedRight\hspace{0pt}\arraybackslash}p{0.102\linewidth}>{\RaggedRight\hspace{0pt}\arraybackslash}p{0.102\linewidth}>{\RaggedRight\hspace{0pt}\arraybackslash}p{0.102\linewidth}>{\RaggedRight\hspace{0pt}\arraybackslash}p{0.102\linewidth}>{\RaggedRight\hspace{0pt}\arraybackslash}p{0.102\linewidth}>{\RaggedRight\hspace{0pt}\arraybackslash}p{0.102\linewidth}>{\RaggedRight\hspace{0pt}\arraybackslash}p{0.102\linewidth}>{\RaggedRight\hspace{0pt}\arraybackslash}p{0.102\linewidth}}
\toprule
Contrast & Series & Pairs & Common contribution & Format contribution & Sum & Frozen Δfull & Maximum discrepancy \\
\midrule
P3-P1 & v2 & 359 & +1.6740 & −3.4395 & −1.7655 & −1.7655 & 1.6e−04 \\
P2-P1 & v2 & 360 & −1.3927 & +4.6739 & +3.2812 & 3.2812 & 1.5e−04 \\
P3-P1 & v1 & 359 & +1.5458 & −3.4577 & −1.9118 & −1.9118 & 1.5e−04 \\
P2-P1 & v1 & 360 & −1.4879 & +4.6800 & +3.1921 & 3.1921 & 1.5e−04 \\
\bottomrule
\end{tabular}
\begin{minipage}{\linewidth}\footnotesize\vspace{2pt}
\textit{Note.} The unit of analysis is a text pair of the contrast. Denominator: 359 or 360 pairs depending on the contrast. Data: series v2 and v1; v1 is given as a stability check of the decomposition, not as the operative result. The decomposition is algebraic and exact: the sum of the contributions matches the frozen contrast to within 1.6·10⁻⁴, which is what the discrepancy column shows. Scale: the sign of a contribution is read under the convention ``higher means more AI-like''. Abbreviations: common — the stylistic feature categories; format — the structural category; Δfull — the full-score contrast from the frozen run. Missing values: none. No correction for multiple comparisons was applied.
\end{minipage}
\end{table*}
